% !TEX root = ../aa_SoM_Chapters.tex

%% HARRIS: Chapter 10

% ö = \%c3\%b6
% ä = \%C3\%A4

\xdef\sectiontitle{\IIsectiontitle} %aiheuttama

%%%%%%%%%%%%%%%
\begin{frame}[label=2_1_a]%[label=2_0_TOC1]
\frametitle{\texorpdfstring{\culine[\sectiontitleulinecolor]{\sectiontitle}}{\sectiontitle} }

\vspace{0.5cm}

\begin{itemize}[leftmargin=0.5cm,itemsep=3mm]
    \item[2.1] \hyperlink{2_1_a}{\CDAlert[blue]{\IIaSubsectiontitle}}
    \item[2.2] \hyperlink{2_2_a}{\IIbSubsectiontitle}	
    \item[2.3] \hyperlink{2_3_a}{\IIcSubsectiontitle}	
    \item[2.4] \hyperlink{2_4_a}{\IIdSubsectiontitle}
    \item[2.5] \hyperlink{2_5_a}{\IIeSubsectiontitle}
    \item[2.6] \hyperlink{2_6_a}{\IIfSubsectiontitle}
    \item[2.7] \hyperlink{2_7_a}{\IIgSubsectiontitle}
    \item[2.8] \hyperlink{2_8_a}{\IIhSubsectiontitle}
    \item[2.9] \hyperlink{2_9_a}{\IIiSubsectiontitle}

\end{itemize}

\endofslide \end{frame}
%%%%%%%%%%%%%%%

\xdef\sectiontitle{\IIaSubsectiontitle} 

%%%%%%%%%%%%%%%
\begin{frame}%[label=2_1_a]
\frametitle{\texorpdfstring{\culine[\sectiontitleulinecolor]{\sectiontitle}}{\sectiontitle} }

\begin{itemize}
	\item Common molecular size is 0.1-1 nm, e.g.:
\begin{itemize}
	\item nitrogen molecule $d=0.37 \mathrm{~nm}$
\item water molecule $d=0.3 \mathrm{~nm}$
\end{itemize}

\item The bond lengths in material are in the range of $\sim 0.1-0.2 \mathrm{~nm}=1-2 \AA$

\item The Bohr radius is $a_{0}=0.0529 \mathrm{~nm}$ (ground state orbit of hydrogen atom)

\item The size of atomic nucleus is $\sim 1-10 \mathrm{fm}=1-10 \times 10^{-6} \mathrm{~nm}$
\end{itemize}
%
%\vspace{2mm}
%%
%\begin{tikzpicture}[remember picture,overlay] % anchor=south
%    \node (A) [yshift=1.45*\figYshift,xshift=\figXshift, anchor=east] at (current page.east) {\includegraphics[page=29,scale=\figscale]{Figures_booklet/Geom/Tikz_main_geom_suom_kolumbus.pdf}};
%\end{tikzpicture}
%

\endofslide 
%\endofsection
\end{frame}
%%%%%%%%%%%%%%%

%%%%%%%%%%%%%%%
\begin{frame}
\frametitle{\texorpdfstring{\culine[\sectiontitleulinecolor]{\sectiontitle}}{\sectiontitle} }

\begin{itemize}
	\item For atoms close to each other, their electrons start seeing a 'multiatom potential energy'
	
	\item Four lowest energy states/wave functions in two neighbouring $L$-sized finite square wells in situations of atoms having varying separation (Fig. 1)
	
\item When coming closer, both wave functions and energy levels are affected. Decreasing the distance, the energy levels will spread further apart
% 6.5 cm = midway, 10 cm = 3/4 
%\vspace{2.5mm}
%\begin{minipage}{4.2cm}
%	levels will spread further apart
%\end{minipage}
%\vspace{3mm}
\end{itemize}

\xdef\loc{south east}
\begin{tikzpicture}[remember picture,overlay] % anchor=south
    \node (A) [yshift=-0.25*\figYshift,xshift=\figXshift, anchor=\loc] at (current page.\loc) {\includegraphics[page=1,width=10.5cm]{Figures_booklet/SOM_10_Bonding_figures/SOM_Bonding_Figure_1.png}}; %scale=\figscale. % 8cm max height
\end{tikzpicture}

\endofslide 
%\endofsection
\end{frame}
%%%%%%%%%%%%%%%

%%%%%%%%%%%%%%%
\begin{frame}
\frametitle{\texorpdfstring{\culine[\sectiontitleulinecolor]{\sectiontitle}}{\sectiontitle} }

\begin{itemize}
	\item If there were three neighbouring finite energy wells, then with smaller separation the case $n=1$ would have three molecular states with\\%[-2mm] 
	% 6.5 cm = midway, 10 cm = 3/4 
\begin{minipage}{7.2cm}
	  slightly different energy levels, which at large separation would converge to three isolated atom $n=1$ states

\item When generalised to a system of $N$ square wells close to each other, then with decreasing separation distance between the square wells, the energy levels will spread and form $N$-state energy bands

\end{minipage}
\vspace{3mm}
\end{itemize}

\xdef\loc{south east}
\begin{tikzpicture}[remember picture,overlay] % anchor=south
    \node (A) [yshift=-0.25*\figYshift,xshift=\figXshift, anchor=\loc] at (current page.\loc) {\includegraphics[page=1,height=7cm]{Figures_booklet/SOM_10_Bonding_figures/SOM_Bonding_Figure_2.png}}; %scale=\figscale. % 8cm max height
\end{tikzpicture}


\endofslide 
%\endofsection
\end{frame}
%%%%%%%%%%%%%%%

%%%%%%%%%%%%%%%
\begin{frame}
\frametitle{\texorpdfstring{\culine[\sectiontitleulinecolor]{\sectiontitle}}{\sectiontitle} }

\begin{itemize}
	% 6.5 cm = midway, 10 cm = 3/4 
\begin{minipage}{8.5cm}
	\item Atoms form molecules when it is energetically favourable, compared to separated atoms. Forming the molecule, valence electrons rearrange to find the lower energy
\end{minipage}
\vspace{3mm}

\item \CDAlert[red]{$\mathrm{H}_{2}{ }^{+}$ molecule}:
\begin{itemize}
	\item A simple example case is $\mathrm{H}_{2}{ }^{+}$ molecule, an ion consisting of two protons and one electron. 
	\item Total energy includes:
\begin{itemize}
	\item Electron energy (kinetic + negative Coulombic attractive potential energy)
\item Coulombic repulsion between two protons
\end{itemize}

\end{itemize}

\end{itemize}

\xdef\loc{north east}
\begin{tikzpicture}[remember picture,overlay] % anchor=south
    \node (A) [yshift=\figYshift,xshift=\figXshift, anchor=\loc] at (current page.\loc) {\includegraphics[page=1,width=5.5cm]{Figures_booklet/SOM_10_Bonding_figures/SOM_Bonding_Figure_3.png}}; %scale=\figscale. % 8cm max height
\end{tikzpicture}

\endofslide 
%\endofsection
\end{frame}
%%%%%%%%%%%%%%%

%%%%%%%%%%%%%%%
\begin{frame}
\frametitle{\texorpdfstring{\culine[\sectiontitleulinecolor]{\sectiontitle}}{\sectiontitle} }

\begin{itemize}
	\item \CDAlert[red]{Electron} in $\mathrm{H}_{2}{ }^{+}$ molecule:
	\item[] \begin{itemize} \small
	\item If the separation between the two protons would be large, then the electron would be with either of the protons, having a Coulombic potential energy of $-13.6~\mathrm{eV}$ 
\item In the other extreme, if the protons 'fused', such as is the case in He atom, then the Coulombic potential energy of the electron would be $-54.4~\mathrm{eV}$
\implies{We expect that the electron is at the lowest energy level, which depends on the proton separation, decreasing with the separation}
\end{itemize}


\item \CDAlert[blue]{Protons} in $\mathrm{H}_{2}{ }^{+}$ molecule:
\item[] \begin{itemize} \small
	\item The positive contribution (\Alert[blue]{repulsion}) to the total energy of the molecule increases as the separation decreases, dramatically, when the separation becomes very small
\item When protons are far away, this is zero
\end{itemize}

\item \CDAlert[violet]{Total}
\item[] \begin{itemize} \small
	\item The molecules such as $\mathrm{H}_{2}{ }^{+}$ exist, thus, apparently there is a minimum potential energy as a function of proton separation
\end{itemize}


\end{itemize}

%\xdef\loc{east}
%\begin{tikzpicture}[remember picture,overlay] % anchor=south
%    \node (A) [yshift=\figYshift,xshift=\figXshift, anchor=\loc] at (current page.\loc) {\includegraphics[page=1,height=7cm]{Figures_booklet/SOM_10_Bonding_figures/SOM_Bonding_Figure_1.png}}; %scale=\figscale. % 8cm max height
%\end{tikzpicture}

\endofslide 
%\endofsection
\end{frame}
%%%%%%%%%%%%%%%

%%%%%%%%%%%%%%%
\begin{frame}
\frametitle{\texorpdfstring{\culine[\sectiontitleulinecolor]{\sectiontitle}}{\sectiontitle} }

\begin{itemize}
	% 6.5 cm = midway, 10 cm = 3/4 
\begin{minipage}{8.5cm}
	\item We know the potential energy of an electron in the Coulombic field of two protons:
$$
U(\mathbf{r})=\frac{1}{4 \pi \varepsilon_{0}} \textrm{((}\frac{-e^{2}}{\left|\mathbf{r}-\frac{1}{2} a \hat{\mathbf{x}}\right|}+\frac{-e^{2}}{\left|\mathbf{r}+\frac{1}{2} a \hat{\mathbf{x}}\right|} \textrm{))}
$$
When this potential energy is plugged
\end{minipage}
\vspace{-2.3mm}
\item[]  into \CDAlert[red]{Schrödinger's equation}, unfortunately we cannot solve it analytically

\item \CDAlert[red]{Numerical solution} reveals that the minimum potential energy for the electron (as a function of the proton separation) will be \CDAlert[red]{$-16.3 \mathrm{eV}$}, and it is obtained \CDAlert[red]{at separation of $a_{\min }=0.11 \mathrm{~nm}$}, which is only about \CDAlert[tunired]{twice the Bohr radius}
\end{itemize}

\xdef\loc{north east}
\begin{tikzpicture}[remember picture,overlay] % anchor=south
    \node (A) [yshift=\figYshift,xshift=\figXshift, anchor=\loc] at (current page.\loc) {\includegraphics[page=1,width=5.5cm]{Figures_booklet/SOM_10_Bonding_figures/SOM_Bonding_Figure_3.png}}; %scale=\figscale. % 8cm max height
\end{tikzpicture}

\endofslide 
%\endofsection
\end{frame}
%%%%%%%%%%%%%%%

%%%%%%%%%%%%%%%
\begin{frame}
\frametitle{\texorpdfstring{\culine[\sectiontitleulinecolor]{\sectiontitle}}{\sectiontitle} }

\begin{itemize}
	\item There will be two 1s-orbital electrons with wave functions $\Psi_{1 s} \textrm{((}r_{1} \textrm{))}$ and $\Psi_{1 s} \textrm{((}r_{2} \textrm{))}$

\item When the protons are closer, we need to take superposition of the two wave functions:\\[3mm]
% 6.5 cm = midway, 10 cm = 3/4 
\begin{minipage}{6.5cm}

	Schrödinger wave equation gives two solutions, \CDAlert[red]{symmetric} \CDAlert[red]{(even)}:
$$
\Psi_{\text {symmetric}}=\frac{1}{\sqrt{2}} \textrm{((}\Psi_{100} \textrm{((}r_{1} \textrm{))}+\Psi_{100} \textrm{((}r_{2} \textrm{))} \textrm{))}
$$
and \CDAlert[blue]{antisymmetric (odd)}:
$$
\Psi_{\text {antisymmetric}}=\frac{1}{\sqrt{2}} \textrm{((}\Psi_{100} \textrm{((}r_{1} \textrm{))}-\Psi_{100} \textrm{((}r_{2} \textrm{))} \textrm{))}
$$
\end{minipage}
\vspace{3mm}
\end{itemize}

\xdef\loc{south east}
\begin{tikzpicture}[remember picture,overlay] % anchor=south
    \node (A) [yshift=-0.25*\figYshift,xshift=\figXshift, anchor=\loc] at (current page.\loc) {\includegraphics[page=1,height=6cm]{Figures_booklet/SOM_10_Bonding_figures/SOM_Bonding_Figure_4.png}}; %scale=\figscale. % 8cm max height
\end{tikzpicture}

\endofslide 
%\endofsection
\end{frame}
%%%%%%%%%%%%%%%

%%%%%%%%%%%%%%%
\begin{frame}
\frametitle{\texorpdfstring{\culine[\sectiontitleulinecolor]{\sectiontitle}}{\sectiontitle} }

\begin{itemize}
	\item The wave function associated with the \CDAlert[red]{ground state} of the molecule \CDAlert[red]{is symmetric}, contributing a majority of its probability/charge density in the region between the protons

\item The attraction which both protons share with the intervening electron cloud, is the root of the molecular bond's lower energy
\end{itemize}

%\xdef\loc{east}
%\begin{tikzpicture}[remember picture,overlay] % anchor=south
%    \node (A) [yshift=\figYshift,xshift=\figXshift, anchor=\loc] at (current page.\loc) {\includegraphics[page=1,height=7cm]{Figures_booklet/SOM_10_Bonding_figures/SOM_Bonding_Figure_1.png}}; %scale=\figscale. % 8cm max height
%\end{tikzpicture}

\endofslide 
%\endofsection
\end{frame}
%%%%%%%%%%%%%%%

%%%%%%%%%%%%%%%
\begin{frame}
\frametitle{\texorpdfstring{\culine[\sectiontitleulinecolor]{\sectiontitle}}{\sectiontitle} }

\begin{itemize}
	\item The construction of the $\mathrm{H}_{2}{ }^{+}$-ion can be extended by adding another electron, thereby \CDAlert[red]{creating a neutral $\mathrm{H}_{2}$ molecule} 
	
	\item Then, there would be two electrons, and in the ground state, both of them would be in the same lowest energy spatial state \implies{ \CDAlert[red]{This energy minimization by sharing of a pair of valence} \CDAlert[red]{electrons of opposite spins} in known as \CDAlert[red]{covalent bond} }

\item The electron pair's spatial state is centrally located and attracts both positive ions (here protons). Both atoms (protons) in $\mathrm{H}_{2}$ 'lay claim' on both electrons, creating a very stable structure

\item The \CDAlert[blue]{same happens with valence $-1$ elements}, similar bonds with two shared electrons are formed in $\mathrm{F}_{2}$, $\mathrm{Cl}_{2}$, $\mathrm{Br}_{2}$ and $\mathrm{I}_{2}$ molecules

\item In $\mathrm{O}_{2}$, the atoms share 4 electrons, in $\mathrm{N}_{2}$, 6 electrons are shared
\end{itemize}

%\xdef\loc{east}
%\begin{tikzpicture}[remember picture,overlay] % anchor=south
%    \node (A) [yshift=\figYshift,xshift=\figXshift, anchor=\loc] at (current page.\loc) {\includegraphics[page=1,height=7cm]{Figures_booklet/SOM_10_Bonding_figures/SOM_Bonding_Figure_1.png}}; %scale=\figscale. % 8cm max height
%\end{tikzpicture}

\endofslide 
%\endofsection
\end{frame}
%%%%%%%%%%%%%%%

%%%%%%%%%%%%%%%
\begin{frame}
\frametitle{\texorpdfstring{\culine[\sectiontitleulinecolor]{\sectiontitle}}{\sectiontitle} }

\begin{itemize}
	\item Certain atoms form noble gas like structures by sharing pairs of valence electrons. However, this is slightly simplified 
	picture, since the wave functions look different from those of noble gases.
	 % 6.5 cm = midway, 10 cm = 3/4 
\item[] \begin{minipage}{6.2cm} 
\item Also, the other, 'nonbonding electrons' may also be shared

\item Recall the wavefunctions $\Psi$ of the $\mathrm{H}_{2}{ }^{+}$-ion described above 

\item The even one has lower energy, since both protons are close to bulk of the electron cloud in the center. It is known as \CDAlert[red]{sigma bonding orbital}, $\sigma 1 s$. %The odd one is known as sigma antibonding orbital, $\sigma^{*} 1 s$. Its energy is higher because the wave function node keeps the electron away from the center region.
\end{minipage}
\vspace{3mm}
\end{itemize}

\xdef\loc{south east}
\begin{tikzpicture}[remember picture,overlay] % anchor=south
    \node (A) [yshift=-0.25*\figXshift,xshift=\figXshift, anchor=\loc] at (current page.\loc) {\includegraphics[page=1,height=6.7cm]{Figures_booklet/SOM_10_Bonding_figures/SOM_Bonding_Figure_4.png}}; %scale=\figscale. % 8cm max height
\end{tikzpicture}

\endofslide 
%\endofsection
\end{frame}
%%%%%%%%%%%%%%%

%%%%%%%%%%%%%%%
\begin{frame}
\frametitle{\texorpdfstring{\culine[\sectiontitleulinecolor]{\sectiontitle}}{\sectiontitle} }

\begin{itemize}
	\item Certain atoms form noble gas like structures by sharing pairs of valence electrons. However, this is slightly simplified 
	picture, since the wave functions look different from those of noble gases.
	 % 6.5 cm = midway, 10 cm = 3/4 
\item[] \begin{minipage}{6.2cm} 
 \item Wave functions look different from those of noble gases. Also, the other, 'nonbonding electrons' may also be shared.

\item The odd one is known as \Alert[blue]{sigma antibonding orbital}, $\sigma^{*} 1 s$. Its energy is higher because the wave function node keeps the electron away from the center region
\end{minipage}
\vspace{3mm}
\end{itemize}

\xdef\loc{south east}
\begin{tikzpicture}[remember picture,overlay] % anchor=south
    \node (A) [yshift=-0.25*\figXshift,xshift=\figXshift, anchor=\loc] at (current page.\loc) {\includegraphics[page=1,height=6.7cm]{Figures_booklet/SOM_10_Bonding_figures/SOM_Bonding_Figure_4.png}}; %scale=\figscale. % 8cm max height
\end{tikzpicture}

\endofslide 
%\endofsection
\end{frame}
%%%%%%%%%%%%%%%

%%%%%%%%%%%%%%%
\begin{frame}
\frametitle{\texorpdfstring{\culine[\sectiontitleulinecolor]{\sectiontitle}}{\sectiontitle} }

\begin{itemize}[itemsep=1.7mm]
% 6.5 cm = midway, 10 cm = 3/4 
\begin{minipage}{10.5cm}
	\item Figure 5 compares the orbitals of the atomic states with both $1 s$ bonding and $1 s$ antibonding states

\item Here the $\pm$-signs do not indicate charge, but the sign of the wave function (which will disappear if one plots probability)

\item The \CDAlert[red]{even $\sigma 1 s$} results from combining \CDAlert[red]{in phase} 

\item The \CDAlert[blue]{odd $\sigma^{*} 1 s$} results from combining a \CDAlert[blue]{half} \CDAlert[blue]{cycle out of phase}

\item The energy diagram of the states show both cases, when the protons are far away from each other and when they are are close by. Behaviour is same as for the $N$ square wells closing up 
\end{minipage}
\vspace{2mm}
\item \CDAlert[red]{Energies spread out with decreasing atomic separation}

\end{itemize}

\xdef\loc{north east}
\begin{tikzpicture}[remember picture,overlay] % anchor=south
    \node (A) [yshift=0.5*\figYshift,xshift=\figXshift, anchor=\loc] at (current page.\loc) {\includegraphics[page=1,width=4.0cm]{Figures_booklet/SOM_10_Bonding_figures/SOM_Bonding_Figure_5.png}}; %scale=\figscale. % 8cm max height
\end{tikzpicture}

\xdef\loc{south east}
\begin{tikzpicture}[remember picture,overlay] % anchor=south
    \node (A) [yshift=-1.25*\figYshift,xshift=\figXshift, anchor=\loc] at (current page.\loc) {\includegraphics[page=1,width=4.0cm]{Figures_booklet/SOM_10_Bonding_figures/SOM_Bonding_Figure_6.png}}; %scale=\figscale. % 8cm max height
\end{tikzpicture}

\endofslide 
%\endofsection
\end{frame}
%%%%%%%%%%%%%%%

%%%%%%%%%%%%%%%
\begin{frame}
\frametitle{\texorpdfstring{\culine[\sectiontitleulinecolor]{\sectiontitle}}{\sectiontitle} }

\begin{itemize}
	\item It is noteworthy to plot the energy levels, along with spins, for the molecules discussed above: $\mathrm{H}_{2}{ }^{+}$, $\mathrm{H}_{2}$, $\mathrm{He}_{2}{ }^{+}$ and $\mathrm{He}_{2}$ (which actually does not exist, since He does not form a molecule)

\item The antibonding orbital has its energy higher than the atomic $1 s$ level

\item $\mathrm{He}_{2}{ }^{+}$-ion can form a bound state, since the two electrons in the $\sigma 1 s$ state and one on the $\sigma^{*} 1 s$ state give a lower total energy than three electrons in atomic $1 s$ states
\end{itemize}

\xdef\loc{south}
\begin{tikzpicture}[remember picture,overlay] % anchor=south
    \node (A) [yshift=-0.25*\figYshift,xshift=0*\figXshift, anchor=\loc] at (current page.\loc) {\includegraphics[page=1,width=11cm]{Figures_booklet/SOM_10_Bonding_figures/SOM_Bonding_Figure_7.png}}; %scale=\figscale. % 8cm max height
\end{tikzpicture}

\endofslide 
%\endofsection
\end{frame}
%%%%%%%%%%%%%%%

%%%%%%%%%%%%%%%
\begin{frame}
\frametitle{\texorpdfstring{\culine[\sectiontitleulinecolor]{\sectiontitle}}{\sectiontitle} }

\begin{itemize}
	% 6.5 cm = midway, 10 cm = 3/4 
\begin{minipage}{9.5cm}
	\item Spherical atomic $2 s$ states give similar behaviour as the $1 s$ states 
	
	\item The $2 p$ states result in more complex behaviour%, as we will see next

\item The three atomic $2 p$ states, perpendicular with each other, are shown in Fig.~8 

\item The angular dependencies of the three equal-energy spatial states of hydrogen are:
$$
\psi_{2,1,0} \propto \cos \theta \qquad and \qquad \psi_{2,1, \pm 1} \propto \sin \theta e^{\pm i \phi}
$$
\end{minipage}
%\vspace{3mm}

\item which may be combined to three equivalent states:
$$
\begin{aligned}
\psi_{2 p_{z}} &=\psi_{2,1,0}  &\propto& ~\cos \theta  \\
 \psi_{2 p_{x}} &= \psi_{2,1,+1}+\psi_{2,1,-1}  &\propto& ~\sin \theta \cos \phi \\ 
 \psi_{2 p_{y}}&=\psi_{2,1,+1}-\psi_{2,1,-1} &\propto& ~\sin \theta \sin \phi
\end{aligned}
$$

\end{itemize}

\xdef\loc{north east}
\begin{tikzpicture}[remember picture,overlay] % anchor=south
    \node (A) [yshift=\figYshift,xshift=\figXshift, anchor=\loc] at (current page.\loc) {\includegraphics[page=1,width=5cm]{Figures_booklet/SOM_10_Bonding_figures/SOM_Bonding_Figure_8.png}}; %scale=\figscale. % 8cm max height
\end{tikzpicture}

\endofslide 
%\endofsection
\end{frame}
%%%%%%%%%%%%%%%

%%%%%%%%%%%%%%%
\begin{frame}
\frametitle{\texorpdfstring{\culine[\sectiontitleulinecolor]{\sectiontitle}}{\sectiontitle} }

% 6.5 cm = midway, 10 cm = 3/4 
\begin{minipage}{10cm}
	\begin{itemize}
	\item For two identical atoms there are six different $2 p$ spatial states

\item The states $2 p_{x}$, having the charge density largest on the molecular axis, are know as $\sigma$-bonds

\item There are both bonding $\sigma 2 p_{x}$ and antibonding $\sigma^{*} 2 p_{x}$ states

\item The two other coordinates have also their $2 p_{y}$ and $2 p_{z}$  states. These are known as \CDAlert[red]{$\pi$-bonds} 

\item These states are: $\pi 2 p_{y}$, $\pi 2 p_{z}$, $\pi^{*} 2 p_{y}$ and $\pi^{*} 2 p_{z}$
\end{itemize}

\end{minipage}
\vspace{3mm}
\xdef\loc{east}
\begin{tikzpicture}[remember picture,overlay] % anchor=south
    \node (A) [yshift=0*\figYshift,xshift=\figXshift, anchor=\loc] at (current page.\loc) {\includegraphics[page=1,height=7cm]{Figures_booklet/SOM_10_Bonding_figures/SOM_Bonding_Figure_9.png}}; %scale=\figscale. % 8cm max height
\end{tikzpicture}

\endofslide 
%\endofsection
\end{frame}
%%%%%%%%%%%%%%%

%%%%%%%%%%%%%%%
\begin{frame}
\frametitle{\texorpdfstring{\culine[\sectiontitleulinecolor]{\sectiontitle}}{\sectiontitle} }

\begin{itemize}
	\item Interestingly, when different pairs of identical two atoms form this kind of covalent bonding, the energy levels of $\sigma$ and $\pi$ states do not necessarily turn out to be in the same order in energy scale for different molecules\\[2mm]

% 6.5 cm = midway, 10 cm = 3/4 
\begin{minipage}{7.25cm}
	\item E.\;g., for $\mathrm{N}_{2}, \mathrm{O}_{2}$ and $\mathrm{F}_{2}$:
\begin{itemize}[itemsep=2mm]

\item $1 s$ electrons are confined to their respective atoms

\item The ordering of levels vary, i.e. $\sigma 2 p$ is higher than $\pi 2 p$ for $\mathrm{N}_{2}$, but lower for $\mathrm{O}_{2}$ and $\mathrm{F}_{2}$

\item The two $\pi^{*} 2 p$ atoms of $\mathrm{O}_{2}$ occupy different states
\end{itemize}
\end{minipage}
\vspace{3mm}

\end{itemize}

\xdef\loc{south east}
\begin{tikzpicture}[remember picture,overlay] % anchor=south
    \node (A) [yshift=-0.25*\figYshift,xshift=\figXshift, anchor=\loc] at (current page.\loc) {\includegraphics[page=1,width=7.5cm]{Figures_booklet/SOM_10_Bonding_figures/SOM_Bonding_Figure_10.png}}; %scale=\figscale. % 8cm max height
\end{tikzpicture}

\endofslide 
%\endofsection
\end{frame}
%%%%%%%%%%%%%%%

%%%%%%%%%%%%%%%
\begin{frame}
\frametitle{\texorpdfstring{\culine[\sectiontitleulinecolor]{\sectiontitle}}{\sectiontitle} }

\begin{itemize}
	\item If the elements of a diatomic molecule are different, then the electrons are not shared equally. This results in \CDAlert[red]{polar covalent bonds}\\[2mm]

% 6.5 cm = midway, 10 cm = 3/4 
\begin{minipage}{9.5cm}
	\item E.g. in $\mathrm{HF}$, the hydrogen's $\mathrm{n}=1$ and fluorine's $\mathrm{n}=2$ electrons are shared asymmetrically

\item The asymmetry is very pronounced in the case where one atom is one electron short, and one atom one too many
\implies{ atoms do not share electrons, but the one atom with a lone electron will 'give' one electron to the other atom } 
\implies{ This results in \CDAlert[blue]{ionic bond} }
%In practice most molecular bonds have their characteristics between these two extremes, between purely covalent and purely ionic bond

\end{minipage}
\vspace{3mm}\end{itemize}

\xdef\loc{east}
\begin{tikzpicture}[remember picture,overlay] % anchor=south
    \node (A) [yshift=\figYshift,xshift=\figXshift, anchor=\loc] at (current page.\loc) {\includegraphics[page=1,width=4.8cm]{Figures_booklet/SOM_10_Bonding_figures/SOM_Bonding_Figure_11.png}}; %scale=\figscale. % 8cm max height
\end{tikzpicture}

\endofslide 
%\endofsection
\end{frame}
%%%%%%%%%%%%%%%

%%%%%%%%%%%%%%%
\begin{frame}
\frametitle{\texorpdfstring{\culine[\sectiontitleulinecolor]{\sectiontitle}}{\sectiontitle} }

\begin{itemize}
	\item So, two different kind of bonds can be formed: \CDAlert[blue]{ionic bonds}  and  \CDAlert[red]{polar covalent bonds}\\[2mm]

% 6.5 cm = midway, 10 cm = 3/4 
\begin{minipage}{9.5cm}
	\item In practice most molecular bonds have their characteristics between these two extremes, between purely covalent and purely ionic bond
\end{minipage}
\vspace{3mm}
\end{itemize}

\xdef\loc{east}
\begin{tikzpicture}[remember picture,overlay] % anchor=south
    \node (A) [yshift=\figYshift,xshift=\figXshift, anchor=\loc] at (current page.\loc) {\includegraphics[page=1,width=4.8cm]{Figures_booklet/SOM_10_Bonding_figures/SOM_Bonding_Figure_11.png}}; %scale=\figscale. % 8cm max height
\end{tikzpicture}

\endofslide 
%\endofsection
\end{frame}
%%%%%%%%%%%%%%%

%%%%%%%%%%%%%%%
\begin{frame}
\frametitle{\texorpdfstring{\culine[\sectiontitleulinecolor]{\sectiontitle}}{\sectiontitle} }

% 6.5 cm = midway, 10 cm = 3/4 
\begin{minipage}{9.5cm}
	\begin{itemize}
	\item Carbon has four $\mathrm{n}=2$ electrons, When surrounded by electron-hungry atoms, carbon shares its four electrons. Spatially this leads to a tetreahedrical configuration of the orbitals, four lobes in four directions.

\item These four states are known as hybrid $s p^{3}$ states. One good example is diamond, which has a very strong crystalline structure\\[-3mm]
$$
{\small%
\begin{aligned}
&\psi_{1}=\psi_{2 s}-\psi_{2 p_{z}} \quad \psi_{11}=\psi_{2 s}+\frac{1}{3} \psi_{2 p_{8}}-\frac{\sqrt{8}}{3} \psi_{2 p_{x}} \\
&\psi_{\mathrm{III}}=\psi_{2 s}+\frac{1}{3} \psi_{2 p_{z}}+\frac{\sqrt{2}}{3} \psi_{2 p_{x}}-\frac{\sqrt{6}}{3} \psi_{2 p_{y}} \\
&\psi_{\mathrm{IV}}=\psi_{2 s}+\frac{1}{3} \psi_{2 p_{z}}+\frac{\sqrt{2}}{3} \psi_{2 p_{x}}+\frac{\sqrt{6}}{3} \psi_{2 p_{y}}
\end{aligned}
}
$$
\end{itemize}
\end{minipage}
%\vspace{3mm}

\xdef\loc{north east}
\begin{tikzpicture}[remember picture,overlay] % anchor=south
    \node (A) [yshift=0.5*\figYshift,xshift=\figXshift, anchor=\loc] at (current page.\loc) {\includegraphics[page=1,width=4.1cm]{Figures_booklet/SOM_10_Bonding_figures/SOM_Bonding_Figure_12.png}}; %scale=\figscale. % 8cm max height
\end{tikzpicture}

\xdef\loc{south east}
\begin{tikzpicture}[remember picture,overlay] % anchor=south
    \node (A) [yshift=-0.25*\figYshift,xshift=\figXshift, anchor=\loc] at (current page.\loc) {\includegraphics[page=1,width=4.1cm]{Figures_booklet/SOM_10_Bonding_figures/SOM_Bonding_Figure_13.png}}; %scale=\figscale. % 8cm max height
\end{tikzpicture}

\endofslide 
%\endofsection
\end{frame}
%%%%%%%%%%%%%%%

%%%%%%%%%%%%%%%
\begin{frame}
\frametitle{\texorpdfstring{\culine[\sectiontitleulinecolor]{\sectiontitle}}{\sectiontitle} }

\begin{itemize}
	\item \CDAlert[red]{Elements near carbon in the periodic system have a tendency} \CDAlert[red]{to form similar geometries}

\item For instance, nitrogen has five $\mathrm{n}=2$ electrons. In ammonia molecule $ \textrm{((}\mathrm{NH}_{3} \textrm{))}$ one of these five takes the place of a 'missing' hydrogen electron. With the proton also missing, the lobe has a negative charge. This gives ammonia an \CDAlert[blue]{electric dipole moment} and accounts for many of its properties as a \CDAlert[tuniblue]{solvent}\\[3mm]

% 6.5 cm = midway, 10 cm = 3/4 
\begin{minipage}{7.5cm}
	\item In \CDAlert[blue]{water}, two extra $\mathrm{n}=2$ electrons of oxygen replace hydrogen's electrons. Two lobes are formed creating also an electric dipole moment
\end{minipage}
\vspace{3mm}
\end{itemize}

\xdef\loc{south east}
\begin{tikzpicture}[remember picture,overlay] % anchor=south
    \node (A) [yshift=-0.25*\figYshift,xshift=\figXshift, anchor=\loc] at (current page.\loc) {\includegraphics[page=1,width=7cm]{Figures_booklet/SOM_10_Bonding_figures/SOM_Bonding_Figure_15.png}}; %scale=\figscale. % 8cm max height
\end{tikzpicture}

%\endofslide 
\endofsection
\end{frame}
%%%%%%%%%%%%%%%
